% from ega4-1-fr.tex
\oldpage[IV-1]{237}

qui n'utilise que (I, 9.4.7). Dans (I, 9.6.6), l'hypothèse «~$X$ est un
schéma quasi-compact~» peut être remplacée par «~$X$ est un préschéma
quasi-compact et quasi-séparé~».
\end{env}

\begin{env}[1.7.10]
\label{IV.1.7.10-fr}
Dans (II, 1.7.15), on peut supposer seulement que $Y$ est quasi-compact et
quasi-séparé, la démonstration n'utilisant que (I, 9.6.5).
\end{env}

\begin{env}[1.7.11]
\label{IV.1.7.11-fr}
Dans la seconde assertion de (II, 3.4.5), on peut substituer aux deux
hypothèses sur $Y$ l'hypothèse plus faible que $Y$ est quasi-compact et
quasi-séparé. De même dans (II, 3.4.8).
\end{env}

\begin{env}[1.7.12]
\label{IV.1.7.12-fr}
Dans (II, 3.8.5), on peut de même supposer seulement que $Y$ est quasi-compact
et quasi-séparé, cette hypothèse intervenant par (I, 9.4.7).
\end{env}

\begin{env}[1.7.13]
\label{IV.1.7.13-fr}
Dans (II, 4.4.1, 4.4.6 et 4.4.7), il suffit de supposer que $Y$ est
\emph{quasi-compact et quasi-séparé}, compte tenu de (1.2.3) dans la
démonstration de (II, 4.4.7). De même, dans (II, 4.4.10.1), il suffit de
supposer $Z$ quasi-compact et quasi-séparé.
\end{env}

\begin{env}[1.7.14]
\label{IV.1.7.14-fr}
Dans (II, 4.5.2) et (II, 4.5.5), il suffit de supposer $X$
\emph{quasi-compact et quasi-séparé}, cette hypothèse intervenant par (I,
9.3.1 et 9.3.2).
\end{env}

\begin{env}[1.7.15]
\label{IV.1.7.15-fr}
Dans (II, 4.6.8), il suffit encore de supposer $X$ \emph{quasi-compact et
quasi-séparé}, cette hypothèse intervenant par (I, 9.4.7).
\end{env}

\begin{env}[1.7.16]
\label{IV.1.7.16-fr}
Dans (II, 5.1.2), il suffit de supposer que $X$ est quasi-compact et
quasi-séparé (l'hypothèse intervenant par (II, 4.5.2 et 4.5.5)). De même dans
(II, 5.1.9), où on utilise (I, 9.6.5).
\end{env}

\begin{env}[1.7.17]
\label{IV.1.7.17-fr}
Dans (II, 5.2.1), il suffit toujours de supposer $X$ quasi-compact et
quasi-séparé (l'hypothèse intervenant par (II, 4.5.2)).
\end{env}

\begin{env}[1.7.18]
\label{IV.1.7.18-fr}
Dans (II, 5.2.2), il faut supposer $f$ quasi-compact et $X$ quasi-séparé~; la
démonstration actuelle est en effet insuffisante (avec les hypothèses faites),
car du fait que pour tout $\mathcal O_X$-Module quasi-cohérent $\mathcal F$ on
a $\mathrm R^1f_*(\mathcal F)=0$, il ne s'ensuit pas nécessairement, pour un
ouvert affine $U$ de $Y$, que la même propriété soit valable pour les
$(\mathcal O_X|f^{-1}(U))$-Modules quasi-cohérents, si on ne sait pas qu'un tel
Module est la \emph{restriction} d'un $\mathcal O_X$-Module quasi-cohérent (la
même remarque s'appliquant aux conditions b) et c'))~; lorsque $f$ est
quasi-compact et $X$ quasi-séparé, ce prolongement est possible en vertu de
(I, 9.4.7), modifié plus haut (1.7.7).
\end{env}

\begin{env}[1.7.19]
\label{IV.1.7.19-fr}
Dans (II, 5.3.2), il suffit de supposer $Y$ \emph{quasi-compact et
quasi-séparé} (ce qui intervient par (II, 4.4.6 et 4.4.7)). De même dans (II,
5.5.3, (ii)).
\end{env}

\begin{env}[1.7.20]
\label{IV.1.7.20-fr}
Dans (III, 1.2.2, 1.2.3 et 1.2.4), on peut remplacer les deux hypothèses sur
$X$ par l'hypothèse moins forte que $X$ est un préschéma \emph{quasi-compact
et quasi-séparé}, la démonstration n'utilisant que (I, 9.3.3).
\end{env}

\begin{env}[1.7.21]
\label{IV.1.7.21-fr}
Dans (III, 1.4.10 à 1.4.14), on peut remplacer «~séparé~» par
«~quasi-séparé~», cette hypothèse servant uniquement à pouvoir appliquer (I,
9.2.2) (voir ci-dessus (1.7.4)). De même, dans (III, 1.4.15), il suffit de
supposer que le morphisme $u$ est \emph{quasi-séparé et quasi-compact}.
\end{env}

\begin{env}[1.7.22]
\label{IV.1.7.22-fr}
Dans (III, 2.1.3), on peut encore remplacer les hypothèses par l'hypothèse
moins forte que $X$ est un préschéma quasi-compact et quasi-séparé, le
raisonnement étant le même que dans (III, 1.2.2).
\end{env}

\subsection{Morphismes de présentation finie et ensembles constructibles}
\label{subsection:IV.1.8-fr}

